\section{结论}
\label{sec:conclusion}

本文系统阐述了纵列双发正交单轴矢量推力模块化飞行器构型的理论基础，按从推进系统物理到
整机控制律的逻辑链展开。主要工作包括：

\begin{enumerate}[label=(\arabic*)]
    \item \textbf{推进系统建模}：从桨盘动量理论和叶素理论出发，建立了推力和反扭矩的
          集总参数模型（$T = k_T\omega^2$，$\tau = k_Q\omega^2 + J_p\dot{\omega}$），
          阐述了量纲分析的相似律（$k_T \propto \rho D^4$，$k_Q \propto \rho D^5$），
          并分析了电机一阶动态和旋转部件陀螺效应的物理机制；
    \item \textbf{推力矢量映射}：通过摆座运动学分析，导出了完整的六维力/力矩映射
          （式\ref{eq:six_component}），分离了主控通道与耦合通道，
          说明推力线过质心和正交摆轴是形成名义主通道结构的重要几何条件，
          但不保证有限摆角下严格解耦；
    \item \textbf{六自由度动力学}：基于牛顿-欧拉方程和四元数姿态运动学，
          建立了整机刚体动力学模型，给出了子步调度的数值积分策略；
    \item \textbf{配平与数值验证}：三方程简化模型得到$\alpha_0=4.123^\circ$、
          $\delta_{t0}=-12.667^\circ$、$\omega_0=268.6$\,rad/s和油门约23.4\%；
          四方程求解得到$\alpha_0=4.1220^\circ$、油门$0.192272969$、
          $\delta_{t0}=-18.7348^\circ$、$\delta_{f0}=-1.5731^\circ$。
          两者约束不同，且四方程尾摆超出固件$\pm15^\circ$限幅。
          Web core本次通过91项自动测试；测试锁定软件行为，但不证明物理参数正确。
          三方程局部正极点$+0.0176657$表示慢发散趋势，20\,s开环空速变化为$+2.234\%$；
    \item \textbf{控制效能分析}：给出了含反扭矩交叉项的完整一阶效能矩阵
          （式\ref{eq:effectiveness_matrix}），并将正交摆轴形成的对角主通道近似与完整模型明确区分；
    \item \textbf{稳定性分析框架}：给出了小扰动线性化方法，解释了纵向两模态（短周期、
          长周期）和横航向三模态（滚转收敛、荷兰滚、螺旋）的物理本质，
          以及反馈对模态的预期影响，并区分连续Jacobian与离散步进矩阵；
    \item \textbf{控制律设计}：区分了Web前飞SAS与固件主控制器。固件使用
          $\bm q^{-1}\otimes\bm q_d$短弧误差，滚转/俯仰四元数外环、偏航ratchet hold、
          deg/s到deg/s$^2$的角速度PI，并在惯量逆解前唯一转换为rad/s$^2$；随后结合
          惯量与陀螺前馈、绝对BTRUE分配、电机模型观测器以及
          $\pm15^\circ/\pm0.7$执行器限幅。
\end{enumerate}

本文的构型贡献在于：\textbf{通过几何设计（正交摆轴、推力线过质心、反向旋转对消）
形成结构清晰的主控制通道，并通过可选机翼模块扩展飞行形态}，使低复杂度直接反馈成为可行的概念方案。
安装机翼时，平台具备固定翼前飞和巡航效率研究的潜力；拆除机翼时，双推进器可承担主要升力，
形成类似多轴飞行器的垂起、悬停与低速机动形态。两种形态共享推进器、摆座和姿态控制基础，
但必须分别标定质量、惯量、推重比和气动参数。
这一设计哲学——用几何换取算法简洁性——在小型电推进飞行器的工程约束下具有实际价值：
其潜在收益是降低名义工作点的控制分配复杂度；是否无需在线优化、能否实现故障隔离，
仍取决于交叉项补偿、饱和管理和失效状态验证。

本文的解析框架不限定具体型号，但数值实现采用\texttt{models/aircraft-model.json}中的
MODEL-DEFAULT参数，全部未经台架或飞行标定，结论应限定为概念级。
文中导出的解析公式（式\ref{eq:six_component}的六维映射、式\ref{eq:effectiveness_matrix}
的效能矩阵、式\ref{eq:coupling_ratio}的耦合比）可作为参数优化、控制律设计
和子系统选型的理论框架。对于具体型号的工程实现，还需在参数标定（台架试验）、
气动数据获取（CFD或风洞）、实时仿真验证（HIL）和飞行试验等多个维度持续迭代。
在控制律方面，INDI、LQR、ADRC与MPC章节保留了方法推导和本平台接口分析；
然而现有对比脚本与固件在四元数误差方向、绝对/增量分配、执行器动态、观测器、调度与限幅上
并不一致，早期LQR/INDI/级联表格及“两个数量级改善”不能作为最终性能结论。
现有首飞与手动直通记录支持控制链已部署及执行器方向正确，但尚无足够的原始日志、重复样本和
参数快照报告超调、调节时间、稳定裕度或完整包线。因此本文当前完成的是构型机理、动力学公式、
固件算法与证据边界的可追溯统一，而不是飞行性能定型。
